% GENERATED FILE. Edit canonical lessons, then run npm run generate:books.
% canonical-source-hash: fnv1a64:8a72715a1945c769
% canonical-lessons: HI-C533-kusti, HI-C533-pahalvan, HI-C533-akhara, HI-C533-larana, HI-C533-paitra

\chapter{In the Wrestling Pit}
\label{ch:hi-in-the-wrestling-pit}
\hlchaptermodality{533}

\begin{chapteropening}
\textbf{By the end of this chapter:} \emph{I can describe a wrestling match.}

Complete the last lesson of chapter 533: I can describe a wrestling match.
\end{chapteropening}

\section[\texorpdfstring{\hi{कुश्ती} (kuśtī)}{कुश्ती (kuśtī)}]{\hi{कुश्ती} (kuśtī) — wrestling; a bout}

\label{lesson:HI-C533-kusti}



\begin{quote}
Before the new one: say the Hindi for a gathering (for music or poetry), then the Hindi for to make dance.
\end{quote}

\subsection*{You'll want to know: \hi{कुश्ती}}
\textbf{\hi{कुश्ती}} — \emph{kuśtī} — \textquotedblleft{}wrestling; a bout\textquotedblright{}.

\begin{quote}
\hi{गाँव} \hi{के} \hi{मेले} \hi{में} \hi{कुश्ती} \hi{होती} \hi{है।}
\end{quote}

\begin{quote}
\emph{Gā̃v ke mele meṁ kuśtī hotī hai.}
\end{quote}

\begin{quote}
There is wrestling at the village fair.
\end{quote}

Feminine.

\begin{cousinweb}
\textbf{\hi{कुश्ती}} (\emph{kuśtī}) is usually traced to Persian \emph{kushtī}, \textquotedblleft{}wrestling\textquotedblright{}.
\end{cousinweb}

\begin{grammarlens}[title={What you've built}]
One more word for describing what you lived and felt.
\end{grammarlens}

\subsection*{Your turn}
\begin{itemize}
\raggedright
  \item \emph{Say it:} \emph{kuśtī}
  \item \emph{Say it:} \emph{kuśtī}, once more
  \item \emph{Say it:} \emph{Gā̃v ke mele meṁ kuśtī hotī hai.} — There is wrestling at the village fair.
  \item \emph{Recall:} say the Hindi for a gathering (for music or poetry), then the Hindi for to make dance, then say \emph{kuśtī} again
\end{itemize}

\begin{tcolorbox}[breakable,colback=teal!4,colframe=teal!35!black,title={Before you move on}]
What does \hi{कुश्ती} mean? (\textquotedblleft{}Wrestling; a bout\textquotedblright{}.) Say it once more.
\end{tcolorbox}
\section[\texorpdfstring{\hi{पहलवान} (pahalvān)}{पहलवान (pahalvān)}]{\hi{पहलवान} (pahalvān) — a wrestler}

\label{lesson:HI-C533-pahalvan}



\begin{quote}
Before the new one: say the Hindi for to make dance, then the Hindi for wrestling.
\end{quote}

\subsection*{You'll want to know: \hi{पहलवान}}
\textbf{\hi{पहलवान}} — \emph{pahalvān} — \textquotedblleft{}a wrestler\textquotedblright{}.

\begin{quote}
\hi{पहलवान} \hi{सुबह} \hi{चार} \hi{बजे} \hi{उठकर} \hi{दौड़ता} \hi{है।}
\end{quote}

\begin{quote}
\emph{Pahalvān subah cār baje uṭhkar dauṛtā hai.}
\end{quote}

\begin{quote}
The wrestler gets up at four in the morning and runs.
\end{quote}

Masculine.

\begin{cousinweb}
\textbf{\hi{पहलवान}} (\emph{pahalvān}) is usually traced to Persian \emph{pahlavān}, \textquotedblleft{}a champion, a hero\textquotedblright{}.
\end{cousinweb}

\begin{grammarlens}[title={What you've built}]
One more word for describing what you lived and felt.
\end{grammarlens}

\subsection*{Your turn}
\begin{itemize}
\raggedright
  \item \emph{Say it:} \emph{pahalvān}
  \item \emph{Say it:} \emph{pahalvān}, once more
  \item \emph{Say it:} \emph{Pahalvān subah cār baje uṭhkar dauṛtā hai.} — The wrestler gets up at four in the morning and runs.
  \item \emph{Recall:} say the Hindi for to make dance, then the Hindi for wrestling, then say \emph{pahalvān} again
\end{itemize}

\begin{tcolorbox}[breakable,colback=teal!4,colframe=teal!35!black,title={Before you move on}]
What does \hi{पहलवान} mean? (\textquotedblleft{}A wrestler\textquotedblright{}.) Say it once more.
\end{tcolorbox}
\section[\texorpdfstring{\hi{अखाड़ा} (akhāṛā)}{अखाड़ा (akhāṛā)}]{\hi{अखाड़ा} (akhāṛā) — a wrestling ground}

\label{lesson:HI-C533-akhara}



\begin{quote}
Before the new one: say the Hindi for wrestling, then the Hindi for a wrestler.
\end{quote}

\subsection*{You'll want to know: \hi{अखाड़ा}}
\textbf{\hi{अखाड़ा}} — \emph{akhāṛā} — \textquotedblleft{}a wrestling ground\textquotedblright{}.

\begin{quote}
\hi{अखाड़े} \hi{की} \hi{मिट्टी} \hi{नरम} \hi{है।}
\end{quote}

\begin{quote}
\emph{Akhāṛe kī miṭṭī naram hai.}
\end{quote}

\begin{quote}
The earth of the wrestling ground is soft.
\end{quote}

Masculine.

\begin{cousinweb}
\textbf{\hi{अखाड़ा}} (\emph{akhāṛā}) is a tadbhava, inherited from Sanskrit \emph{akṣavāṭa}, \textquotedblleft{}a wrestling ground\textquotedblright{} (CDIAL 39).
\end{cousinweb}

\begin{grammarlens}[title={What you've built}]
One more word for describing what you lived and felt.
\end{grammarlens}

\subsection*{Your turn}
\begin{itemize}
\raggedright
  \item \emph{Say it:} \emph{akhāṛā}
  \item \emph{Say it:} \emph{akhāṛā}, once more
  \item \emph{Say it:} \emph{Akhāṛe kī miṭṭī naram hai.} — The earth of the wrestling ground is soft.
  \item \emph{Recall:} say the Hindi for wrestling, then the Hindi for a wrestler, then say \emph{akhāṛā} again
\end{itemize}

\begin{tcolorbox}[breakable,colback=teal!4,colframe=teal!35!black,title={Before you move on}]
What does \hi{अखाड़ा} mean? (\textquotedblleft{}A wrestling ground\textquotedblright{}.) Say it once more.
\end{tcolorbox}
\section[\texorpdfstring{\hi{लड़ाना} (laṛānā)}{लड़ाना (laṛānā)}]{\hi{लड़ाना} (laṛānā) — to pit against; to make fight}

\label{lesson:HI-C533-larana}



\begin{quote}
Before the new one: say the Hindi for a wrestler, then the Hindi for a wrestling ground.
\end{quote}

\subsection*{You'll want to know: \hi{लड़ाना}}
\textbf{\hi{लड़ाना}} — \emph{laṛānā} — \textquotedblleft{}to pit against; to make fight\textquotedblright{}.

\begin{quote}
\hi{बच्चों} \hi{ने} \hi{छत} \hi{पर} \hi{दो} \hi{पतंगें} \hi{लड़ाईं।}
\end{quote}

\begin{quote}
\emph{Baccoṁ ne chat par do pataṅgeṁ laṛāīṁ.}
\end{quote}

\begin{quote}
On the roof the children pitted two kites against each other.
\end{quote}

Past, with \hi{ने} (he or she …): \hi{उसने} \hi{दो} \hi{पहलवानों} \hi{को} \hi{लड़ाया}, \emph{usne do pahalvānoṁ ko laṛāyā}.

\begin{cousinweb}
\textbf{\hi{लड़ाना}} (\emph{laṛānā}) is the causative of \emph{laṛnā}, \textquotedblleft{}to fight\textquotedblright{}, from \emph{*laḍ}, \textquotedblleft{}fight\textquotedblright{}, a form CDIAL reconstructs from its descendants (CDIAL 10920).
\end{cousinweb}

\begin{grammarlens}[title={What you've built}]
One more word for describing what you lived and felt.
\end{grammarlens}

\subsection*{Your turn}
\begin{itemize}
\raggedright
  \item \emph{Say it:} \emph{laṛānā}
  \item \emph{Say it:} \emph{laṛānā}, once more
  \item \emph{Say it:} \emph{Baccoṁ ne chat par do pataṅgeṁ laṛāīṁ.} — On the roof the children pitted two kites against each other.
  \item \emph{Recall:} say the Hindi for a wrestler, then the Hindi for a wrestling ground, then say \emph{laṛānā} again
\end{itemize}

\begin{tcolorbox}[breakable,colback=teal!4,colframe=teal!35!black,title={Before you move on}]
What does \hi{लड़ाना} mean? (\textquotedblleft{}To pit against; to make fight\textquotedblright{}.) Say it once more.
\end{tcolorbox}
\section[\texorpdfstring{\hi{पैंतरा} (paĩtrā)}{पैंतरा (paĩtrā)}]{\hi{पैंतरा} (paĩtrā) — a stance; a move (in wrestling)}

\label{lesson:HI-C533-paitra}



\begin{quote}
Before the new one: say the Hindi for a wrestling ground, then the Hindi for to pit against.
\end{quote}

\subsection*{You'll want to know: \hi{पैंतरा}}
\textbf{\hi{पैंतरा}} — \emph{paĩtrā} — \textquotedblleft{}a stance; a move (in wrestling)\textquotedblright{}.

\begin{quote}
\hi{पहलवान} \hi{ने} \hi{एक} \hi{नया} \hi{पैंतरा} \hi{दिखाया।}
\end{quote}

\begin{quote}
\emph{Pahalvān ne ek nayā paĩtrā dikhāyā.}
\end{quote}

\begin{quote}
The wrestler showed off a new move.
\end{quote}

Masculine.

\begin{cousinweb}
\textbf{\hi{पैंतरा}} (\emph{paĩtrā}) is a tadbhava, inherited from Sanskrit \emph{pādāntara}, \textquotedblleft{}the space of one step\textquotedblright{} (CDIAL 8073). CDIAL gives the Hindi word as \emph{paitrā}, \textquotedblleft{}a way of placing the feet when turning in fencing or wrestling\textquotedblright{}.
\end{cousinweb}

\begin{grammarlens}[title={What you've built}]
One more word for describing what you lived and felt.
\end{grammarlens}

\subsection*{Your turn}
\begin{itemize}
\raggedright
  \item \emph{Say it:} \emph{paĩtrā}
  \item \emph{Say it:} \emph{paĩtrā}, once more
  \item \emph{Say it:} \emph{Pahalvān ne ek nayā paĩtrā dikhāyā.} — The wrestler showed off a new move.
  \item \emph{Recall:} say the Hindi for a wrestling ground, then the Hindi for to pit against, then say \emph{paĩtrā} again
\end{itemize}

\begin{tcolorbox}[breakable,colback=teal!4,colframe=teal!35!black,title={Before you move on}]
What does \hi{पैंतरा} mean? (\textquotedblleft{}A stance; a move (in wrestling)\textquotedblright{}.) Say it once more.
\end{tcolorbox}
